\documentclass[12pt,a4paper]{report}
\usepackage[utf8]{inputenc}
\usepackage[french]{babel}
\usepackage[T1]{fontenc}
\usepackage{geometry}
\geometry{top=2.5cm, bottom=2.5cm, left=2.5cm, right=2.5cm}
\usepackage{graphicx}
\usepackage{xcolor}
\usepackage{listings}
\usepackage{hyperref}
\usepackage{fancyhdr}
\usepackage{titlesec}
\usepackage{booktabs}
\usepackage{tcolorbox}

% Configuration des couleurs
\definecolor{primary}{HTML}{1A365D}   % Bleu foncé
\definecolor{secondary}{HTML}{2B6CB0} % Bleu clair
\definecolor{accent}{HTML}{319795}    % Teal
\definecolor{bgcode}{HTML}{F7FAFC}    % Gris très clair
\definecolor{commentgreen}{HTML}{38A169}
\definecolor{keywordblue}{HTML}{3182CE}

% Configuration hyperref
\hypersetup{
    colorlinks=true,
    linkcolor=primary,
    filecolor=accent,
    urlcolor=secondary,
    pdftitle={Rapport Projet de Fin de Module - dApp TP3},
}

% Configuration du code source (Listings)
\lstdefinelanguage{Solidity}{
    keywords={contract, abstract, library, is, struct, mapping, address, uint, uint256, int, string, bool, function, returns, public, private, external, internal, view, pure, payable, memory, storage, override, virtual, require, constructor, bytes, pragma, solidity, SPDX-License-Identifier},
    keywordstyle=\color{keywordblue}\bfseries,
    ndkeywords={Rectangle, Forme, GestionChaines, Exercice1\_Addition, Exercice2\_Conversion, Exercice4\_EstPositif, Exercice5\_Parite, Exercice6\_Tableau, Payment},
    ndkeywordstyle=\color{accent}\bfseries,
    identifierstyle=\color{black},
    sensitive=true,
    comment=[l]{//},
    morecomment=[s]{/*}{*/},
    commentstyle=\color{commentgreen}\ttfamily,
    stringstyle=\color{red}\ttfamily,
    morestring=[b]",
    morestring=[b]'
}

\lstset{
    backgroundcolor=\color{bgcode},
    basicstyle=\ttfamily\small,
    breakatwhitespace=false,
    breaklines=true,
    captionpos=b,
    keepspaces=true,
    numbers=left,
    numberstyle=\tiny\color{gray},
    numbersep=5pt,
    showspaces=false,
    showstringspaces=false,
    showtabs=false,
    tabsize=2,
    frame=single,
    rulecolor=\color{gray!30}
}

% En-tête et pied de page
\pagestyle{fancy}
\fancyhf{}
\fancyhead[L]{\textcolor{primary}{\bfseries Projet dApp TP3}}
\fancyhead[R]{\textcolor{secondary}{Solidity \& React}}
\fancyfoot[C]{\thepage}
\renewcommand{\headrulewidth}{0.4pt}
\renewcommand{\footrulewidth}{0.4pt}

% Styling des titres de section
\titleformat{\section}
  {\normalfont\large\bfseries\color{secondary}}
  {\thesection}{1em}{}

% Page de garde
\begin{document}

\begin{titlepage}
    \centering
    \vspace*{1.5cm}
    {\Huge\bfseries\color{primary} Projet de Fin de Module}\\[0.5cm]
    {\Large\bfseries\color{secondary} Développement d'une Application Décentralisée (dApp)}\\[1cm]
    \rule{\linewidth}{1.5pt}\\[1.5cm]

    {\large\bfseries Réalisé dans le cadre du :}\\[0.3cm]
    {\large TP3 -- Smart Contracts Solidity \& Intégration Web3}\\[1.5cm]

    \begin{tcolorbox}[colback=primary!5,colframe=primary,arc=5pt,title=Technologies Utilisées,halign=center]
        \textbf{Smart Contracts :} Solidity (v0.8.20) \\
        \textbf{Développement \& Tests :} Truffle Suite \\
        \textbf{Blockchain Locale :} Ganache (port 7545) \\
        \textbf{Interface Utilisateur :} React (Vite, Web3.js)
    \end{tcolorbox}

    \vfill
    \rule{\linewidth}{1pt}\\[0.5cm]
    \begin{minipage}{0.45\textwidth}
        \begin{flushleft} \large
            \textbf{Préparé par :}\\
            [Étudiant / Développeur]
        \end{flushleft}
    \end{minipage}
    \hfill
    \begin{minipage}{0.45\textwidth}
        \begin{flushright} \large
            \textbf{Encadré par :}\\
            M. OUALLA
        \end{flushright}
    \end{minipage}\\[1cm]

    {\large \today}
\end{titlepage}

\tableofcontents
\newpage

% ============================================================
\chapter{Introduction et Architecture}
% ============================================================

\section{Introduction}
L'objectif de ce projet est de concevoir et développer une application décentralisée (dApp) complète
regroupant un ensemble d'exercices pratiques (8 exercices au total) dans le cadre du TP3 du module
\textit{Blockchain et Web3} (Filière CI-ILIA, S4, AU 2025/2026).
Chaque exercice permet d'explorer des concepts fondamentaux des Smart Contracts Solidity et de leur
intégration avec une interface web moderne développée en React.

\section{Architecture du Projet}
Le projet suit une structure moderne séparant clairement la logique blockchain
(contrats intelligents, compilateur Truffle, réseau local Ganache) de la logique client
(React, Vite, Web3.js).

\begin{itemize}
    \item \textbf{contracts/ :} Contient les 8 contrats Solidity (\texttt{.sol}).
    \item \textbf{migrations/ :} Contient les 8 scripts de déploiement Truffle.
    \item \textbf{frontend/ :} Application React/Vite avec 8 pages d'exercice et une page d'accueil.
    \item \textbf{update-env.js :} Script Node.js qui lit les adresses déployées depuis \texttt{build/contracts/*.json}
          et met à jour automatiquement \texttt{frontend/.env}.
    \item \textbf{truffle-config.js :} Configuration du compilateur Solidity (v0.8.20) et du réseau Ganache.
\end{itemize}

\section{Procédure de Démarrage}
Le flux de travail pour démarrer l'application locale se compose des étapes suivantes :
\begin{enumerate}
    \item \textbf{Ganache :} Démarrage d'une blockchain Ethereum locale sur le port \texttt{7545}.
    \item \textbf{Compilation \& Déploiement :} Exécution de \texttt{truffle compile} puis \texttt{truffle migrate}.
    \item \textbf{Mise à jour des adresses :} Exécution de \texttt{node update-env.js} pour synchroniser le frontend.
    \item \textbf{Lancement de l'interface client :} Exécution de \texttt{npm run dev} dans le dossier \texttt{frontend}.
\end{enumerate}

\newpage

% ============================================================
\chapter{Détail des Exercices}
% ============================================================

% -----------------------------------------------------------
\section{Exercice 1 : Somme de deux variables}
% -----------------------------------------------------------
L'objectif est de créer un contrat avec deux nombres comme variables d'état et de déclarer
deux fonctions pour calculer la somme.

\begin{itemize}
    \item \textbf{\texttt{addition1()} -- view :} Calcule la somme des deux variables d'état
          \texttt{a} et \texttt{b} déclarées dans le contrat. Elle utilise le modificateur
          \texttt{view} car elle lit l'état sans le modifier.
    \item \textbf{\texttt{addition2(uint x, uint y)} -- pure :} Prend deux nombres en paramètres
          et retourne leur somme, sans accéder ni modifier l'état du contrat (\texttt{pure}).
    \item \textbf{\texttt{setValues(uint \_a, uint \_b)} :} Fonction utilitaire permettant de
          modifier dynamiquement les variables d'état depuis le frontend.
\end{itemize}

\subsection{Code du Smart Contract -- \texttt{Exercice1\_Addition.sol}}
\begin{lstlisting}[language=Solidity, caption={Contrat Addition.sol -- fonctions view et pure}]
// SPDX-License-Identifier: MIT
pragma solidity ^0.8.20;

// Exercice 1 : Somme de deux nombres
contract Exercice1_Addition {
    uint public a = 10;
    uint public b = 20;

    // Retourne la somme des deux variables d'etat (view)
    function addition1() public view returns (uint) {
        return a + b;
    }

    // Retourne la somme de deux parametres (pure)
    function addition2(uint x, uint y) public pure returns (uint) {
        return x + y;
    }

    // Permet de modifier les variables d'etat
    function setValues(uint _a, uint _b) public {
        a = _a;
        b = _b;
    }
}
\end{lstlisting}

\subsection{Capture d'écran de la page Exercice 1}
\begin{figure}[h!]
    \centering
    % \includegraphics[width=0.85\textwidth]{images/exercice1.png}
    \begin{tcolorbox}[colback=red!5,colframe=red!50!black,title=Capture d'écran manquante]
        Insérez ici la capture d'écran de l'Exercice 1 (\texttt{images/exercice1.png})
    \end{tcolorbox}
    \caption{Interface de calcul de somme -- Exercice 1}
\end{figure}

\newpage

% -----------------------------------------------------------
\section{Exercice 2 : Conversion Ether \& Wei}
% -----------------------------------------------------------
Cet exercice porte sur la conversion des unités de mesure de la blockchain Ethereum,
entre l'Ether et le Wei ($1\ \text{Ether} = 10^{18}\ \text{Wei}$). Le contrat ne contient
\textbf{aucune variable d'état} ; toutes ses fonctions sont \texttt{pure}.

\begin{itemize}
    \item \textbf{\texttt{etherEnWei(uint montantEther)} :} Multiplie le montant en Ether par
          \texttt{1 ether} (constante Solidity $= 10^{18}$) et retourne l'équivalent en Wei.
    \item \textbf{\texttt{weiEnEther(uint montantWei)} :} Effectue la conversion inverse en divisant
          par \texttt{1 ether}.
\end{itemize}

\subsection{Code du Smart Contract -- \texttt{Exercice2\_Conversion.sol}}
\begin{lstlisting}[language=Solidity, caption={Contrat Conversion.sol -- Ether et Wei}]
// SPDX-License-Identifier: MIT
pragma solidity ^0.8.20;

// Exercice 2 : Conversion de cryptomonnaies
contract Exercice2_Conversion {

    // Convertit des Ethers en Wei (montantEther * 1 ether = montantEther * 10^18)
    function etherEnWei(uint montantEther) public pure returns (uint) {
        return montantEther * 1 ether;
    }

    // Convertit des Wei en Ether (division par 10^18)
    function weiEnEther(uint montantWei) public pure returns (uint) {
        return montantWei / 1 ether;
    }
}
\end{lstlisting}

\subsection{Capture d'écran de la page Exercice 2}
\begin{figure}[h!]
    \centering
    % \includegraphics[width=0.85\textwidth]{images/exercice2.png}
    \begin{tcolorbox}[colback=red!5,colframe=red!50!black,title=Capture d'écran manquante]
        Insérez ici la capture d'écran de l'Exercice 2 (\texttt{images/exercice2.png})
    \end{tcolorbox}
    \caption{Convertisseur instantané Ether $\leftrightarrow$ Wei -- Exercice 2}
\end{figure}

\newpage

% -----------------------------------------------------------
\section{Exercice 3 : Gestion de chaînes de caractères}
% -----------------------------------------------------------
L'objectif est de se familiariser avec la gestion des chaînes (\texttt{string}) en Solidity.
Le contrat \texttt{GestionChaines} expose les fonctions suivantes :

\begin{itemize}
    \item \textbf{\texttt{setMessage()} / \texttt{getMessage()} :} Écriture et lecture de la
          variable d'état \texttt{message}.
    \item \textbf{\texttt{concatener(a, b)} :} Concatène deux chaînes passées en paramètres
          grâce à \texttt{string.concat(a, b)}.
    \item \textbf{\texttt{concatenerAvec(autre)} :} Concatène \texttt{message} avec une autre chaîne.
    \item \textbf{\texttt{longueur(s)} :} Retourne la longueur via \texttt{bytes(s).length}.
    \item \textbf{\texttt{comparer(a, b)} :} Compare deux chaînes via \texttt{keccak256(abi.encodePacked(...))}
          et retourne un \texttt{bool}.
\end{itemize}

\subsection{Code du Smart Contract -- \texttt{Exercice3\_GestionChaines.sol}}
\begin{lstlisting}[language=Solidity, caption={Contrat GestionChaines.sol -- manipulation de string}]
// SPDX-License-Identifier: MIT
pragma solidity ^0.8.20;

// Exercice 3 : Gestion des chaines de caracteres
contract GestionChaines {
    string public message;

    // Modifier la valeur de message
    function setMessage(string memory _message) public {
        message = _message;
    }

    // Retourner la valeur de message
    function getMessage() public view returns (string memory) {
        return message;
    }

    // Concatener deux chaines passees en parametres
    function concatener(string memory a, string memory b)
        public pure returns (string memory) {
        return string.concat(a, b);
    }

    // Concatener message avec une autre chaine
    function concatenerAvec(string memory autre)
        public view returns (string memory) {
        return string.concat(message, autre);
    }

    // Retourner la longueur d'une chaine (bytes(s).length)
    function longueur(string memory s) public pure returns (uint) {
        return bytes(s).length;
    }

    // Comparer deux chaines via keccak256
    function comparer(string memory a, string memory b)
        public pure returns (bool) {
        return keccak256(abi.encodePacked(a))
               == keccak256(abi.encodePacked(b));
    }
}
\end{lstlisting}

\subsection{Capture d'écran de la page Exercice 3}
\begin{figure}[h!]
    \centering
    % \includegraphics[width=0.85\textwidth]{images/exercice3.png}
    \begin{tcolorbox}[colback=red!5,colframe=red!50!black,title=Capture d'écran manquante]
        Insérez ici la capture d'écran de l'Exercice 3 (\texttt{images/exercice3.png})
    \end{tcolorbox}
    \caption{Interface de gestion des chaînes de caractères -- Exercice 3}
\end{figure}

\newpage

% -----------------------------------------------------------
\section{Exercice 4 : Signe d'un nombre (Est Positif)}
% -----------------------------------------------------------
Ce contrat définit une unique fonction \texttt{estPositif()} qui prend un entier relatif
(\texttt{int}) en paramètre et retourne \texttt{true} si ce nombre est supérieur ou égal à zéro,
et \texttt{false} sinon.

\subsection{Code du Smart Contract -- \texttt{Exercice4\_EstPositif.sol}}
\begin{lstlisting}[language=Solidity, caption={Contrat EstPositif.sol -- vérification du signe}]
// SPDX-License-Identifier: MIT
pragma solidity ^0.8.20;

// Exercice 4 : Verifier si un nombre est positif
contract Exercice4_EstPositif {

    // Retourne true si n >= 0, false sinon
    function estPositif(int n) public pure returns (bool) {
        return n >= 0;
    }
}
\end{lstlisting}

\subsection{Capture d'écran de la page Exercice 4}
\begin{figure}[h!]
    \centering
    % \includegraphics[width=0.85\textwidth]{images/exercice4.png}
    \begin{tcolorbox}[colback=red!5,colframe=red!50!black,title=Capture d'écran manquante]
        Insérez ici la capture d'écran de l'Exercice 4 (\texttt{images/exercice4.png})
    \end{tcolorbox}
    \caption{Interface d'évaluation du signe d'un nombre -- Exercice 4}
\end{figure}

\newpage

% -----------------------------------------------------------
\section{Exercice 5 : Parité d'un nombre (Pair / Impair)}
% -----------------------------------------------------------
Utilisation de l'opérateur modulo (\texttt{\%}) pour déterminer si un entier naturel
(\texttt{uint}) est pair ou impair. Deux fonctions booléennes sont exposées.

\subsection{Code du Smart Contract -- \texttt{Exercice5\_Parite.sol}}
\begin{lstlisting}[language=Solidity, caption={Contrat Parite.sol -- opérateur modulo}]
// SPDX-License-Identifier: MIT
pragma solidity ^0.8.20;

// Exercice 5 : Verifier la parite d'un nombre entier
contract Exercice5_Parite {

    // Retourne true si le nombre est pair
    function estPair(uint n) public pure returns (bool) {
        return n % 2 == 0;
    }

    // Retourne true si le nombre est impair
    function estImpair(uint n) public pure returns (bool) {
        return n % 2 != 0;
    }
}
\end{lstlisting}

\subsection{Capture d'écran de la page Exercice 5}
\begin{figure}[h!]
    \centering
    % \includegraphics[width=0.85\textwidth]{images/exercice5.png}
    \begin{tcolorbox}[colback=red!5,colframe=red!50!black,title=Capture d'écran manquante]
        Insérez ici la capture d'écran de l'Exercice 5 (\texttt{images/exercice5.png})
    \end{tcolorbox}
    \caption{Interface de test de parité -- Exercice 5}
\end{figure}

\newpage

% -----------------------------------------------------------
\section{Exercice 6 : Gestion des Tableaux dynamiques}
% -----------------------------------------------------------
Ce contrat stocke une liste de nombres dans un tableau dynamique \texttt{uint[] nombres}
et expose les opérations suivantes :

\begin{itemize}
    \item Le constructeur initialise le tableau avec les valeurs \{10, 20, 30, 40, 50\}.
    \item \textbf{\texttt{ajouterNombre(uint n)} :} Ajoute un élément au tableau via \texttt{push()}.
    \item \textbf{\texttt{getElement(uint index)} :} Retourne l'élément à l'indice donné ;
          la fonction \texttt{require()} protège contre les indices hors limites.
    \item \textbf{\texttt{afficheTableau()} :} Retourne l'intégralité du tableau.
    \item \textbf{\texttt{calculerSomme()} :} Parcourt le tableau et retourne la somme totale.
    \item \textbf{\texttt{getTaille()} :} Retourne la taille courante du tableau (\texttt{length}).
\end{itemize}

\subsection{Code du Smart Contract -- \texttt{Exercice6\_Tableau.sol}}
\begin{lstlisting}[language=Solidity, caption={Contrat Tableau.sol -- tableaux dynamiques}]
// SPDX-License-Identifier: MIT
pragma solidity ^0.8.20;

// Exercice 6 : Tableau dynamique de nombres
contract Exercice6_Tableau {
    uint[] public nombres;

    // Initialiser le tableau dans le constructeur
    constructor() {
        nombres.push(10);
        nombres.push(20);
        nombres.push(30);
        nombres.push(40);
        nombres.push(50);
    }

    // Ajouter un nombre au tableau
    function ajouterNombre(uint n) public {
        nombres.push(n);
    }

    // Retourner l'element a l'indice index (avec require)
    function getElement(uint index) public view returns (uint) {
        require(index < nombres.length,
                "Erreur : index hors limites du tableau");
        return nombres[index];
    }

    // Retourner tout le tableau
    function afficheTableau() public view returns (uint[] memory) {
        return nombres;
    }

    // Calculer la somme de tous les elements
    function calculerSomme() public view returns (uint) {
        uint somme = 0;
        for (uint i = 0; i < nombres.length; i++) {
            somme += nombres[i];
        }
        return somme;
    }

    // Retourner la taille du tableau
    function getTaille() public view returns (uint) {
        return nombres.length;
    }
}
\end{lstlisting}

\subsection{Capture d'écran de la page Exercice 6}
\begin{figure}[h!]
    \centering
    % \includegraphics[width=0.85\textwidth]{images/exercice6.png}
    \begin{tcolorbox}[colback=red!5,colframe=red!50!black,title=Capture d'écran manquante]
        Insérez ici la capture d'écran de l'Exercice 6 (\texttt{images/exercice6.png})
    \end{tcolorbox}
    \caption{Interface interactive pour tableaux dynamiques -- Exercice 6}
\end{figure}

\newpage

% -----------------------------------------------------------
\section{Exercice 7 : POO, Héritage et Formes Géométriques}
% -----------------------------------------------------------
Cet exercice illustre la \textbf{Programmation Orientée Objet} en Solidity à travers un contrat
abstrait \texttt{Forme} et un contrat concret \texttt{Rectangle} qui en hérite.

\begin{itemize}
    \item \textbf{Contrat abstrait \texttt{Forme} :} Déclare les coordonnées \texttt{x}, \texttt{y},
          un constructeur, \texttt{deplacerForme()}, \texttt{afficheXY()}, une fonction
          \texttt{virtual pure afficheInfos()} retournant \og{}Je suis une forme\fg{},
          et une fonction \texttt{virtual surface()} à implémenter.
    \item \textbf{Contrat \texttt{Rectangle is Forme} :} Ajoute \texttt{lo} (longueur) et
          \texttt{la} (largeur), implémente \texttt{surface()} $=$ \texttt{lo} $\times$ \texttt{la},
          redéfinit \texttt{afficheInfos()} en retournant \og{}Je suis Rectangle\fg{},
          et ajoute \texttt{afficheLoLa()}.
\end{itemize}

\subsection{Code du Smart Contract -- \texttt{Exercice7\_Formes.sol}}
\begin{lstlisting}[language=Solidity, caption={Contrat Formes.sol -- héritage et POO}]
// SPDX-License-Identifier: MIT
pragma solidity ^0.8.20;

// Exercice 7 : POO -- Formes geometriques

// Contrat abstrait Forme
abstract contract Forme {
    uint public x;
    uint public y;

    constructor(uint _x, uint _y) {
        x = _x;
        y = _y;
    }

    // Deplacer la forme
    function deplacerForme(uint dx, uint dy) public {
        x += dx;
        y += dy;
    }

    // Retourner les coordonnees actuelles
    function afficheXY() public view returns (uint, uint) {
        return (x, y);
    }

    // Fonction virtuelle pure -- a redefinir dans les sous-classes
    function afficheInfos() public virtual pure returns (string memory) {
        return "Je suis une forme";
    }

    // Fonction virtuelle -- a implementer dans Rectangle
    function surface() public virtual view returns (uint);
}

// Contrat concret Rectangle heritant de Forme
contract Rectangle is Forme {
    uint public lo; // longueur
    uint public la; // largeur

    constructor(uint _x, uint _y, uint _longueur, uint _largeur)
        Forme(_x, _y)
    {
        lo = _longueur;
        la = _largeur;
    }

    // Calculer la surface du rectangle : lo * la
    function surface() public view override returns (uint) {
        return lo * la;
    }

    // Redefinir afficheInfos
    function afficheInfos() public pure override returns (string memory) {
        return "Je suis Rectangle";
    }

    // Retourner la longueur et la largeur
    function afficheLoLa() public view returns (uint, uint) {
        return (lo, la);
    }
}
\end{lstlisting}

\subsection{Capture d'écran de la page Exercice 7}
\begin{figure}[h!]
    \centering
    % \includegraphics[width=0.85\textwidth]{images/exercice7.png}
    \begin{tcolorbox}[colback=red!5,colframe=red!50!black,title=Capture d'écran manquante]
        Insérez ici la capture d'écran de l'Exercice 7 (\texttt{images/exercice7.png})
    \end{tcolorbox}
    \caption{Interface géométrique interactive -- Exercice 7}
\end{figure}

\newpage

% -----------------------------------------------------------
\section{Exercice 8 : Transactions financières -- \texttt{msg.sender} \& \texttt{msg.value}}
% -----------------------------------------------------------
L'objectif est d'utiliser les variables globales \texttt{msg.sender} et \texttt{msg.value}
pour gérer les transactions Ether. Le contrat \texttt{Payment} fonctionne comme suit :

\begin{itemize}
    \item \textbf{Variable \texttt{recipient} :} Adresse du destinataire initialisée dans le constructeur.
    \item \textbf{\texttt{receivePayment()} payable :} Reçoit des Ethers et vérifie via
          \texttt{require()} que \texttt{msg.value > 0}.
    \item \textbf{\texttt{withdraw()} :} Vérifie que \texttt{msg.sender == recipient} puis
          transfère la totalité du solde du contrat au destinataire via
          \texttt{payable(recipient).transfer(address(this).balance)}.
    \item \textbf{\texttt{getBalance()} :} Retourne le solde courant du contrat en Wei.
\end{itemize}

\subsection{Code du Smart Contract -- \texttt{Exercice8\_Payment.sol}}
\begin{lstlisting}[language=Solidity, caption={Contrat Payment.sol -- msg.sender et msg.value}]
// SPDX-License-Identifier: MIT
pragma solidity ^0.8.20;

// Exercice 8 : Gestion des transactions avec msg.sender et msg.value
contract Payment {
    address public recipient;

    constructor(address _recipient) {
        recipient = _recipient;
    }

    // Recevoir des Ethers -- verifie que msg.value > 0
    function receivePayment() public payable {
        require(msg.value > 0,
                "Erreur : le montant envoye doit etre superieur a 0");
    }

    // Retirer les fonds vers le destinataire
    function withdraw() public {
        require(msg.sender == recipient,
                "Erreur : seul le destinataire peut retirer les fonds");
        payable(recipient).transfer(address(this).balance);
    }

    // Consulter le solde du contrat
    function getBalance() public view returns (uint) {
        return address(this).balance;
    }
}
\end{lstlisting}

\subsection{Capture d'écran de la page Exercice 8}
\begin{figure}[h!]
    \centering
    % \includegraphics[width=0.85\textwidth]{images/exercice8.png}
    \begin{tcolorbox}[colback=red!5,colframe=red!50!black,title=Capture d'écran manquante]
        Insérez ici la capture d'écran de l'Exercice 8 (\texttt{images/exercice8.png})
    \end{tcolorbox}
    \caption{Passerelle de paiement décentralisée -- Exercice 8}
\end{figure}

\newpage

% ============================================================
\chapter{Conclusion}
% ============================================================
Ce projet a permis d'explorer concrètement le processus d'ingénierie logicielle Web3.
Les 8 exercices ont couvert progressivement les bases de Solidity~: modificateurs \texttt{view}/\texttt{pure},
conversions d'unités, manipulation de \texttt{string}, arithmétique booléenne, tableaux dynamiques,
programmation orientée objet avec héritage, et enfin la gestion sécurisée des transactions via
\texttt{msg.sender} et \texttt{msg.value}.

La séparation des couches -- calcul décentralisé d'un côté, rendu dynamique React de l'autre --
illustre la puissance et les défis inhérents au développement sur Ethereum. Les smart contracts
sécurisent l'exécution des règles métiers, tandis que le frontend React (connecté via \texttt{Web3.js})
offre à l'utilisateur final une interface fluide et réactive vers les fonctions du réseau blockchain.

\end{document}
